\documentclass[../cours_spe_trigo.tex]{subfiles}
\begin{document}
\setcounter{chapter}{14}
\chapter{Fonctions réciproques : \texorpdfstring{$\arcsin$, $\arccos$ et $\arctan$}{arcsin, arccos et arctan}}
\section{Pourquoi faut-il restreindre les fonctions ?}
Une fonction admet une fonction réciproque sur un ensemble seulement si elle y est bijective. Or la fonction sinus prend la même valeur en une infinité d'angles ; il faut donc la restreindre à un intervalle où elle est strictement monotone.
\begin{definition}
Les trois fonctions réciproques usuelles sont définies par les choix suivants :
\[
\begin{aligned}
\arcsin &: [-1,1]\longrightarrow[-\pi/2,\pi/2],\\
\arccos &: [-1,1]\longrightarrow[0,\pi],\\
\arctan &: \mathbb R\longrightarrow]-\pi/2,\pi/2[.
\end{aligned}
\]
Sur ces intervalles, chacune renvoie \textbf{l'unique angle} dont le sinus, le cosinus ou la tangente est l'argument.
\end{definition}
\textbf{Exemples.} $\arcsin(1/2)=\pi/6$. En revanche,
\[\arcsin(\sin(5\pi/6))=\pi/6\ne5\pi/6.\]
De même, $\arccos(\cos(5\pi/3))=\pi/3$.
\section{Dérivées des fonctions réciproques (approfondissement)}
\begin{demonstration}
Posons $y=\arcsin x$ pour $-1<x<1$. Alors $x=\sin y$ et $y\in]-\pi/2,\pi/2[$. En dérivant : $1=\cos y\,y'$. Puisque $\cos y>0$ dans cet intervalle et $\cos^2y=1-\sin^2y=1-x^2$, on a $\cos y=\sqrt{1-x^2}$. Ainsi
\[\boxed{(\arcsin x)'=\frac1{\sqrt{1-x^2}}.}\]
On démontre de même, avec $y=\arccos x$ et $\sin y>0$ pour $0<y<\pi$ :
\[\boxed{(\arccos x)'=-\frac1{\sqrt{1-x^2}}.}\]
Enfin, si $y=\arctan x$, $x=\tan y$, donc $1=(1+\tan^2y)y'=(1+x^2)y'$ et
\[\boxed{(\arctan x)'=\frac1{1+x^2}.}\]
$\square$
\end{demonstration}
\end{document}
